% Section 6. Discussion and conclusion.

\section{Discussion}
\label{sec:discussion}

In both the worked example and the three external datasets, each sample failed at least one check tied to a use researchers make of such samples.
In the worked example none of these failures could be seen in a single answer.
Our recommendation for researchers running persona studies is to begin before the generation step, first by establishing whether human data exist to ground the sample, and second by deciding which claims about the population the study intends to make.
Holding the model's ability to stand in for those people in critical doubt, in good faith, the researcher then checks each claim against the rung designed to test it.
Group gaps, population levels, single cases and scale scores are separate extrapolations, and each one needs its own check.

Using the ladder depends on meeting its requirements.
Every level above the gate compares the simulated sample against human data, and without a reference on the same instrument the ladder cannot be applied.
Each rung additionally carries its own requirements, set out in Figure~\ref{fig:ladder}; where a study's data cannot meet them, that rung is not run and the claims it would support are not made.
Because the reference determines which personas can be compared, how many draws each needs and which gaps can be measured precisely enough to certify, it has to be in hand before generation.
Additionally, the researcher fixes how each persona attribute maps to a reference variable before generation and reports that mapping with the verdicts, because in the worked example matching income on the label's dollar amount rather than the poverty-ratio band changed the verdict on one contrast. % R: l2_verdicts.csv
Personas built on real respondents' profiles, as \citet{bisbee2024synthetic} and \citet{argyle2023outofone} did, offer the cleanest reference by giving every simulated answer a human twin.

Our controls and the three external datasets address whether the ladder itself can be trusted.
Pseudo-models built from real respondents rarely fail, while failures planted into those same respondents are caught at the rung each one targets.
Applied to other authors' releases without changing a rule, the ladder ran on instruments, models and designs it was not built around.
In those releases it supports the partisan gaps reported by \citet{bisbee2024synthetic} and \citet{argyle2023outofone}, while most gaps by sex, age and race either fail or cannot be read.
We do not dispute those studies' findings, though the results point toward a narrower interpretation of their claims of fidelity, centred on the partisan comparisons.

In practice the ladder strengthens the context around each claim drawn from a persona sample and weeds out readings the data do not support.
The right value for a gap or a level still has to come from the human data.

Regarding the ladder's own limitations, its thresholds were fixed before the external analyses and act as defaults, to be replaced by a tolerance a user can justify for their own use.
Certifying a pass also takes more draws, more personas per group and a more precise reference than detecting a failure, and an unresolved verdict can reflect the design as well as the simulator.
Under the design of the worked example, with 30 draws per persona and a national survey as reference, level 2 serves mainly to detect distorted gaps.
Finally, the frozen ladder does not read contrasts with no population gap, leaving untested any gap a model invents where people show none; an equivalence test on the simulated gap in raw units would cover this case.

\subsection*{Conclusion}

The validity ladder replaces a single judgement of whether a simulated sample resembles people with a separate check for each use of the sample, each carrying a pass rule, a human reference and the claim it supports.
Across 28,800 PHQ-8 assessments from four models, the ladder read every model's persona means and failed every model on each use above them. % R: corpus_v3_provenance.csv; l1_personfit_main.csv
Applied to three published silicon samples, it narrowed what each release can be used to claim and identified claims some releases cannot test at all.
For a researcher, the result is a list of the claims a sample supports, each tied to the check behind it.
We release the code for every rung and the receipt for every number, allowing grounded or fine-tuned simulators to be checked against the same baseline.
